\section*{\centering\fontsize{14}{21}\selectfont\bfseries KATA PENGANTAR}

Puji dan syukur kami panjatkan ke hadirat Tuhan Yang Maha Esa atas rahmat dan karunia-Nya, sehingga kami dapat menyelesaikan penyusunan proposal ini dengan tepat waktu. Proposal ini disusun dalam rangka mengikuti kompetisi \textit{Business Plan Competition} (BPC) IFEST 2026 yang diselenggarakan oleh Himpunan Mahasiswa Teknik Informatika Universitas Padjadjaran (HIMATIF UNPAD).

Gagasan bisnis ini dirancang sebagai upaya untuk menghadirkan solusi inovatif dalam menghadapi masalah inefisiensi logistik pengangkutan sampah perkotaan dan krisis kapasitas Tempat Pemrosesan Akhir (TPA) Bantar Gebang, sekaligus menciptakan peluang usaha yang memiliki nilai ekonomis, daya saing tinggi, dan berdampak positif bagi masyarakat luas.

Keberhasilan penyusunan proposal ini tidak lepas dari bantuan dan dukungan dari berbagai pihak. Pada kesempatan ini, kami ingin menyampaikan apresiasi dan terima kasih yang sebesar-besarnya kepada:
\begin{enumerate}[leftmargin=1.5cm]
  \item Keluarga besar kami yang senantiasa memberikan dukungan moral maupun material selama proses penyusunan proposal ini. \textit{(--- lengkapi dengan nama dosen pembimbing/pendamping, jika ada, sebelum submit ---)}
  \item Panitia penyelenggara IFEST BPC 2026 HIMATIF UNPAD atas kesempatan dan bimbingan teknis yang diberikan.
  \item Seluruh pihak yang telah memberikan masukan, data, maupun dukungan lain yang tidak dapat kami sebutkan satu per satu.
\end{enumerate}

Kami menyadari bahwa proposal ini masih memiliki keterbatasan dan jauh dari kata sempurna. Oleh karena itu, kami sangat terbuka terhadap kritik, saran, serta masukan konstruktif dari dewan juri dan pembaca sekalian demi penyempurnaan rancangan bisnis ini di masa depan.

Besar harapan kami agar proposal bisnis ini dapat dinilai secara objektif, memberikan manfaat nyata, serta dapat direalisasikan di masa depan.

\vspace{1cm}
\begin{flushright}
Jatinangor, 2026\\
Tim Penyusun
\end{flushright}
